%% annex/p05.tex --- proof annex for Part 5 (pilot plus Phase-4 depth visit)
%% One annexitem per proved=/sketched= body result of Part 5, in body
%% order.  The diagonal lemma is cited and therefore carries commentary
%% rather than an annex item.
\hbannexchapter{5}{Arithmetic and Incompleteness}\label{annex:p05}

The items below follow the order of the Part-5 body chapters.  Sketches
open with a note enumerating their omissions; the diagonal lemma is a
cited result and therefore carries commentary rather than an item.

\begin{annexitem}{annex:p5.1}{thm:p05-q-bounded-cases}{Bounded-case enumeration in \(\mathsf Q\)}
\begin{proof}
Metainduction on \(n\), each instance a \(\mathsf Q\)-derivation from
the axioms \(Sx\neq0\), \(Sx=Sy\to x=y\), \(x\neq0\to\exists y\,x=Sy\),
\(x+0=x\), \(x+Sy=S(x+y)\).  Base: \(z+y=\bar0\) forces \(y=0\), since
\(y=Sw\) gives \(S(z+w)=z+Sw=\bar0\), contradicting \(Sx\neq0\).  Step
to \(n+1\): from \(z+y=S\bar n\), either \(y=0\), or \(y=Sw\) with
\(S(z+w)=S\bar n\); injectivity of \(S\) gives \(z+w=\bar n\), so
\(w\leq\bar n\), the instance for \(n\) lists \(w\), and applying \(S\)
to its disjuncts completes the list for \(y\).
\end{proof}
\end{annexitem}

\begin{annexitem}{annex:p5.2}{thm:p05-q-dichotomy}{Dichotomy in \(\mathsf Q\)}
\begin{proof}
Metainduction on \(n\) from the same axioms.  Base: the witness
\(z\defeq y\) gives \(y+\bar0=y\), so \(\bar0\leq y\).  Step: if
\(y=0\), the witness \(S\bar n\) gives \(y\leq S\bar n\); if \(y=Sw\),
the instance for \(n\) at \(w\) gives either \(z+w=\bar n\), whence
\(z+Sw=S\bar n\) and \(y\leq S\bar n\), or \(u+\bar n=w\), whence
\(u+S\bar n=S(u+\bar n)=y\) and \(S\bar n\leq y\).
\end{proof}
\end{annexitem}

% ---- Tier-3 commentary (no annex item: the diagonal lemma is cited, and
% the bidirectional 112=112 discipline reserves items for proved=/sketched=
% results).  Referenced from the body method box def:p05-diagonal-method.
\phantomsection\label{annex:p5-diagonal-commentary}%
\noindent\emph{Commentary: the construction behind the diagonal lemma
(Theorem~\ref{thm:p05-diagonal-lemma}, cited from
\cite{BoolosBurgessJeffrey2007}).}\enspace
Let \(D(x,y)\) represent the function \(\mathrm{diag}\) in \(\mathsf Q\),
so that for each standard \(n\),
\[
  \mathsf Q\derives\forall y\,
  \bigl(D(\bar n,y)\leftrightarrow
  y=\overline{\mathrm{diag}(n)}\bigr).
\]
Given \(\psi(x)\), define
\(\beta(x)\defeq\exists y\,(D(x,y)\land\psi(y))\), let
\(b=\ulcorner\beta\urcorner\), and let \(G\) be the sentence
\(\beta(\bar b)\); by construction
\(\mathrm{diag}(b)=\ulcorner G\urcorner\).  Representability of
\(\mathrm{diag}\) gives, already in \(\mathsf Q\),
\[
\begin{aligned}
  G&\leftrightarrow\exists y\,(D(\bar b,y)\land\psi(y))\\
   &\leftrightarrow\psi(\overline{\mathrm{diag}(b)})
    \leftrightarrow\psi(\overline{\ulcorner G\urcorner}),
\end{aligned}
\]
and \(T\supseteq\mathsf Q\) proves the same biconditional.  The cited
proof supplies the primitive-recursive verification and numeralwise
representability of \(\mathrm{diag}\) on which this explanation depends.

\begin{annexitem}{annex:p5.3}{thm:p05-goedel-first}{G\"odel's first incompleteness theorem, original form}
\sketchnote{the complexity-preserving fixed point and the
representability theorem enter as cited
(Theorem~\ref{thm:p05-representability},
\cite{Shoenfield1967,BoolosBurgessJeffrey2007}).}
(1) If \(T\derives G_T\), a concrete \(T\)-proof exists, so
representation of the proof relation gives
\(T\derives\mathrm{Prov}_T(\overline{\ulcorner G_T\urcorner})\), while
(G) gives its negation---\(T\) is inconsistent.

(2) If \(T\) is consistent, (1) gives
\(\mathcal N_{\mathrm{ar}}\entails\neg\mathrm{Prov}_T(\overline{\ulcorner
G_T\urcorner})\) (no standard number codes a proof); (G) holds in
\(\mathcal N_{\mathrm{ar}}\) since \(\mathcal N_{\mathrm{ar}}\entails\mathsf Q\), so
\(\mathcal N_{\mathrm{ar}}\entails G_T\) and \(\neg G_T\) is a false \(\Sigma^0_1\)
sentence, unprovable in a \(\Sigma^0_1\)-sound \(T\).

(3) If \(T\derives\neg G_T\), (G) gives
\(T\derives\exists p\,\mathrm{Proof}_T(p,\overline{\ulcorner
G_T\urcorner})\); consistency and (1) give
\(T\derives\neg\mathrm{Proof}_T(\bar n,\overline{\ulcorner
G_T\urcorner})\) for every standard \(n\)---violating omega-consistency.
\end{annexitem}

\begin{annexitem}{annex:p5.4}{thm:p05-rosser}{Rosser's theorem}
\sketchnote{cited inputs: diagonal lemma and representability theorem;
Lemmas~\ref{thm:p05-q-bounded-cases} and~\ref{thm:p05-q-dichotomy}
carry the bounded comparisons.}
By the diagonal lemma, fix \(R_T\)---every proof of \(R_T\) has a proof
of \(\neg R_T\) with no larger code (Convention~\ref{def:p05-leq})---with
\[
  \mathsf Q\derives R_T\leftrightarrow\forall y\,\bigl(
  \mathrm{Proof}_T(y,\overline{\ulcorner R_T\urcorner})\to
  \exists z\leq y\,
  \mathrm{Proof}_T(z,\overline{\ulcorner\neg R_T\urcorner})\bigr).
\]

If \(T\derives R_T\) with least proof code \(p\), representation and
the fixed point give \(T\derives\exists z\leq\bar p\,
\mathrm{Proof}_T(z,\overline{\ulcorner\neg R_T\urcorner})\), while
consistency and representation give
\(T\derives\neg\mathrm{Proof}_T(\bar q,\overline{\ulcorner\neg
R_T\urcorner})\) for each \(q\leq_{\N}p\);
Lemma~\ref{thm:p05-q-bounded-cases} then refutes the bounded
existential disjunct by disjunct---inconsistent, with no metatheoretic
proof of \(\neg R_T\) produced.

If \(T\derives\neg R_T\) with least proof code \(q\), representation
gives \(T\derives\mathrm{Proof}_T(\bar q,\overline{\ulcorner\neg
R_T\urcorner})\).  Inside \(T\), from
\(\mathrm{Proof}_T(y,\overline{\ulcorner R_T\urcorner})\),
Lemma~\ref{thm:p05-q-dichotomy} gives \(y\leq\bar q\lor\bar q\leq y\):
if \(y\leq\bar q\), Lemma~\ref{thm:p05-q-bounded-cases} lists \(y\)
among \(\bar0,\ldots,\bar q\), consistency and representation refute
each case, and the bounded existential follows from the contradiction;
if \(\bar q\leq y\), then \(z\defeq\bar q\) is a witness; discharging,
the fixed point yields \(T\derives R_T\)---again inconsistent.
\end{annexitem}

\begin{annexitem}{annex:p5.5}{thm:p05-theoremhood-undecidable}{Undecidability of theoremhood}
\sketchnote{rests on the Rosser sketch (Annex~\ref{annex:p5.4}); its
bounded proof-comparison verification is the omission.}
The theorem set is computably enumerable by enumerating finite
\(T\)-proofs.  If it were decidable, enumerate the sentences as
\(\sigma_0,\sigma_1,\ldots\); starting from \(D_0=\varnothing\), add
\(\sigma_n\) to \(D_n\) when \(T\cup D_n\cup\{\sigma_n\}\) is
consistent---decidable via the deduction theorem---and \(\neg\sigma_n\)
otherwise; one choice always preserves consistency.  Then
\(T^*=T\cup\bigcup_nD_n\) is a consistent, computably enumerable,
deductively complete extension of \(\mathsf Q\), and Rosser's theorem
applied to \(T^*\) yields an independent sentence---a contradiction.
\end{annexitem}

%% annex/p05/depth-completion.tex --- Phase-4 P5 depth proofs

\begin{annexitem}{annex:p5.models.1}{thm:p05-compactness-nonstandard}{Compactness produces nonstandard models}
\begin{proof}
Adjoin \(c\) and set \(S=T\cup\{c\neq\bar n:n\in\N\}\).  Finite subsets
expand \(\mathcal N_{\mathrm{ar}}\) by choosing \(c\) outside their
numerals, so compactness gives \(\mathcal M'\entails S\).  Its reduct is
nonstandard: an isomorphism to \(\mathcal N_{\mathrm{ar}}\) would map \(c\)
to \(n\) and, preserving numerals, force \(c=\bar n\).  Downward
L\"owenheim--Skolem in the expanded countable signature, with \(c\)
designated, gives a countable elementary submodel whose reduct satisfies
\(T\) and retains \(c\).
\end{proof}
\end{annexitem}

\begin{annexitem}{annex:p5.sigma.1}{thm:p05-sigma-one-completeness}{External and formalized Sigma-one completeness}
\sketchnote{omits the full induction evaluating closed bounded formulas
in \(\mathsf Q\), the chosen normal-form conversion, and the exact
primitive-recursive construction of proof codes from witness tuples,
including its numeral-substitution invariant and verification inside
\(\PA\).}
Write a true closed \(\Sigma^0_1\) sentence as
\(\exists\vec x\,\varphi(\vec x)\) with bounded \(\varphi\), and choose
its standard witnesses.  Closed terms and atomic formulas evaluate on
their numerals; Boolean induction and the bounded-case lemmas of
Section~\ref{sec:p05-diagonalization-prereq} turn the true bounded
instance into a \(\mathsf Q\)-derivation.  Existential introduction gives
the external clause.  For the internal clause, the same structural
construction is arithmetized: from a witness tuple and its bounded
evaluation data it computes a \(\PA\)-proof code for the corresponding
instance, and then appends existential introductions.  PA verifies these
constructors and their composition for the fixed sentence, yielding the
displayed formalized implication.
\end{annexitem}


\begin{annexitem}{annex:p5.6}{thm:p05-lob}{L\"ob's theorem}
\begin{proof}
By the diagonal lemma, choose \(\psi\) with
\[
  T\derives\psi\leftrightarrow
  \bigl(\mathrm{Prov}_T(\ulcorner\psi\urcorner)\to\varphi\bigr).
  \tag{L1}
\]
D1 and D2 applied to the forward direction of (L1) give
\(T\derives\mathrm{Prov}_T(\ulcorner\psi\urcorner)\to
\mathrm{Prov}_T(\ulcorner\mathrm{Prov}_T(\ulcorner\psi\urcorner)
\to\varphi\urcorner)\); combining with D3 by internalized modus ponens
(D2) gives
\(T\derives\mathrm{Prov}_T(\ulcorner\psi\urcorner)\to
\mathrm{Prov}_T(\ulcorner\varphi\urcorner)\).  The hypothesis then
gives \(T\derives\mathrm{Prov}_T(\ulcorner\psi\urcorner)\to\varphi\),
so the converse direction of (L1) yields \(T\derives\psi\); D1 gives
\(T\derives\mathrm{Prov}_T(\ulcorner\psi\urcorner)\), whence
\(T\derives\varphi\).
\end{proof}
\end{annexitem}

\begin{annexitem}{annex:p5.7}{thm:p05-goedel-second}{G\"odel's second incompleteness theorem}
\begin{proof}
Suppose \(T\derives\mathsf{Con}(T)\).  In the fixed explosive calculus,
\(\mathsf{Con}(T)\) propositionally yields
\(T\derives\mathrm{Prov}_T(\overline{\ulcorner\bottom_A\urcorner})
\to\bottom_A\), the hypothesis of Theorem~\ref{thm:p05-lob} with
\(\varphi=\bottom_A\); hence \(T\derives\bottom_A\), contradicting
consistency.
\end{proof}
\end{annexitem}

\begin{annexitem}{annex:p5.8}{thm:p05-tarski}{Tarski's undefinability theorem}
\begin{proof}
If \(\mathrm{Tr}(x)\) were such a formula, the diagonal lemma applied
to \(\neg\mathrm{Tr}(x)\) gives a sentence \(L\) with
\(\mathsf Q\derives L\leftrightarrow
\neg\mathrm{Tr}(\overline{\ulcorner L\urcorner})\), true in
\(\mathcal N_{\mathrm{ar}}\) since \(\mathcal N_{\mathrm{ar}}\entails\mathsf Q\); the assumed truth
definition gives \(\mathcal N_{\mathrm{ar}}\entails
\mathrm{Tr}(\overline{\ulcorner L\urcorner})\leftrightarrow L\).
Combining the two, \(\mathcal N_{\mathrm{ar}}\entails L\leftrightarrow\neg L\),
impossible classically.
\end{proof}
\end{annexitem}

\begin{annexitem}{annex:p5.9}{thm:p05-truth-undefinable}{Undefinability of the arithmetical truth set}
\begin{proof}
A formula defining the displayed set in \(\mathcal N_{\mathrm{ar}}\) would be an
arithmetical truth definition; and computably enumerable sets are
\(\Sigma^0_1\)-definable, so the set is not computably enumerable.
\end{proof}
\end{annexitem}
